\documentclass[11pt]{article}
\usepackage[a4paper,margin=2.5cm]{geometry}
\usepackage{amsmath,amssymb}
\usepackage{graphicx}
\usepackage{booktabs}
\usepackage{parskip}

\title{MA-INF 2201 Computer Vision WS26/27 \\ Exercise Sheet 2 --- Report / Discussion Sections}
\author{Mohamadhosein Sabaghi \and Sepehr Noey Kordkandi}
\date{}

\begin{document}
\maketitle

\noindent\textit{Note: the Parseval's theorem proof (Question 1) is submitted separately as
\texttt{template1.pdf}. This report covers the discussion/explanation prompts from Questions
2--5.}

\setcounter{section}{1} % next \section becomes "2", matching the sheet's Question 2

% ==========================================================================
\section{Fourier Transform and Spatial vs.\ Frequency Filtering}

\subsection{Magnitude and Phase (Part 2a)}

\begin{table}[h]
\centering
\begin{tabular}{lcc}
\toprule
Reconstruction & MAD vs.\ Image 1 & MAD vs.\ Image 2 \\
\midrule
Magnitude(1) + Phase(2) & $61.29$ & $\mathbf{37.06}$ \\
Magnitude(2) + Phase(1) & $\mathbf{35.63}$ & $50.04$ \\
\bottomrule
\end{tabular}
\end{table}

In both cases, the reconstruction is closer (lower MAD, bold) to whichever original image
supplied its \emph{phase}, not whichever supplied its magnitude: ``Magnitude(1)+Phase(2)''
resembles Image 2 more ($37.06 < 61.29$), and ``Magnitude(2)+Phase(1)'' resembles Image 1
more ($35.63 < 50.04$). This confirms the classical result that \textbf{phase}, not
magnitude, carries the recognisable spatial structure of an image --- the magnitude
spectrum mainly encodes overall contrast/texture energy per frequency, while the phase
encodes \emph{where} those frequency components align in the image, which is what the
visual system (and pixel-wise comparison) primarily responds to.

MAD does capture this effect in the correct direction, but only partially: the gap between
the two MAD values for each reconstruction ($61.29$ vs.\ $37.06$, and $35.63$ vs.\ $50.04$)
is modest relative to the overall magnitude of both numbers (out of a $0$--$255$ intensity
range), even though the phase-matched image is typically much more clearly recognisable by
eye than this gap alone would suggest. This is consistent with MAD being a purely
pixel-wise error metric that does not model perceptual/structural similarity --- it reacts
to the right component being present, but understates how large the subjective difference
in recognisability actually is.

\subsection{Spatial vs.\ Frequency Domain Filtering (Part 2b)}

\begin{center}
\begin{tabular}{lc}
\toprule
Filter & MAD (spatial vs.\ frequency domain) \\
\midrule
Box ($9\times9$) & $3.52\times10^{-14}$ \\
Gaussian ($9\times9$, $\sigma=1.5$) & $2.64\times10^{-14}$ \\
\bottomrule
\end{tabular}
\end{center}

Both values are many orders of magnitude below the required threshold of $10^{-7}$, at the
level of floating-point round-off rather than any real discrepancy. This confirms that the
frequency-domain implementation correctly performs \emph{linear} convolution (via explicit
zero-padding to the full convolution size before multiplying spectra, then cropping back to
the original image size) rather than circular convolution, matching the direct spatial-domain
sliding-window convolution to numerical precision.

% ==========================================================================
\section{Gaussian and Laplacian Pyramids}

\subsection{Gaussian Pyramid and Aliasing (Part 3a)}

Gaussian filter and settings used: $5\times5$ kernel, $\sigma=1.0$, reflect boundary padding.
Downsampling: decimation by 2 (keep every other row/column).

Downsampling by a factor of 2 halves the sampling rate, which also halves the Nyquist
frequency the result can represent. A high-frequency synthetic pattern (checkerboard /
stripes) contains energy above this new Nyquist limit. Without pre-filtering, that excess
high-frequency content does not simply disappear --- it folds back (\emph{aliases}) into
lower frequencies, producing a distorted or moir\'e-like pattern rather than a faithfully
shrunk version of the original. Pre-filtering with a Gaussian low-pass suppresses frequency
content above the new Nyquist limit before the samples are discarded, so the filtered
pyramid shows a clean, correctly-shrunk pattern at every level. The natural image
(\texttt{kitty.png}), having comparatively little very-high-frequency content, shows the
same effect much more subtly, mainly as slightly softer edges in the filtered pyramid
versus marginally jagged edges in the unfiltered one.
% TODO: adjust/extend this paragraph once you've visually compared the saved
% gaussian_pyramid_kitty.png and gaussian_pyramid_pattern.png figures.

\subsection{Laplacian Pyramid and Reconstruction (Part 3b)}

Settings used for decomposition and reconstruction (identical in both directions, as
required): Gaussian filter $5\times5$, $\sigma=1.0$, reflect boundary padding; expand
operation via \texttt{cv2.resize} with bilinear interpolation (\texttt{cv2.INTER\_LINEAR}),
resizing directly to each target level's exact shape. No clipping or quantisation was
applied to any intermediate image.

\begin{center}
\textbf{Reconstruction MAD: $2.08 \times 10^{-19}$}
\end{center}

This is effectively zero (floating-point noise level), confirming that expanding and
summing the Laplacian levels from coarse to fine exactly reconstructs the original image,
as expected from the telescoping definition $L_\ell = G_\ell - \text{expand}(G_{\ell+1})$.

% ==========================================================================
\section{Canny Edge Detector}

\textbf{Settings used (identical for both detectors):} low threshold $=50$, high threshold
$=100$; Gaussian smoothing $5\times5$, $\sigma=1.0$, applied once and reused as input for
both our implementation and \texttt{cv2.Canny} (so the two detectors start from exactly the
same preprocessed, \texttt{uint8} image); $3\times3$ Sobel kernels; L2 gradient magnitude
($\texttt{L2gradient=True}$ in the \texttt{cv2.Canny} call, matching our
$\sqrt{d_x^2+d_y^2}$); 8-connected hysteresis propagation from strong to weak edges.

\begin{center}
\begin{tabular}{lc}
\toprule
Metric & Value \\
\midrule
F1-score (target $\geq 0.6$) & $\mathbf{0.6117}$ \\
MAD (target $\leq 0.07$) & $\mathbf{0.0454}$ \\
\bottomrule
\end{tabular}
\end{center}

Both targets are met. An earlier version of our non-maximum suppression step had the $45^\circ$
and $135^\circ$ diagonal neighbour pairs swapped (checking the anti-diagonal neighbours for
a gradient oriented along the main diagonal and vice versa), which mis-thinned diagonally-oriented
edges into fragments without noticeably changing the overall pixel-wise MAD (diagonal edges are a
small fraction of total image pixels). After correcting the neighbour assignment to match the true
gradient direction, F1 improved from $0.4231$ to $0.6117$ and MAD improved from $0.0586$ to
$0.0454$, confirming that accurate directional suppression --- not just correct thresholds --- is
essential for the edge maps to agree with \texttt{cv2.Canny} at the pixel level.

% ==========================================================================
\section{Chamfer Distance}

\subsection{Chamfer 5-7-11 Distance Transform (Part 5a)}

Edges were obtained from \texttt{uni-bonn.jpg} via \texttt{cv2.Canny} (thresholds $100$,
$200$), and the Chamfer 5-7-11 transform was computed with a two-pass raster scan (forward:
top-to-bottom/left-to-right; backward: bottom-to-top/right-to-left), each pass propagating
distances using the axis ($5$), diagonal ($7$), and knight-move ($11$) step costs from the
already-visited neighbourhood. The raw result was divided by $5$ to express it in
approximate pixel units for comparison against \texttt{cv2.distanceTransform} (Euclidean
metric, computed on the inverted edge mask so edge pixels are zeros):

\begin{center}
\begin{tabular}{lc}
\toprule
Metric & Value \\
\midrule
Mean absolute difference (Chamfer/5 vs.\ Euclidean) & $0.0164$ \\
Max absolute difference (Chamfer/5 vs.\ Euclidean) & $0.2387$ \\
\bottomrule
\end{tabular}
\end{center}

The difference is small but not identically zero, as expected: the Chamfer 5-7-11 transform
only \emph{approximates} true Euclidean distance using a small, fixed set of local step costs
along specific directions (axis-aligned, diagonal, knight-move). These costs are chosen to
closely match real Euclidean step lengths ($5:7:11 \approx 5:7.07:11.18$), but any shortest
path to the nearest edge that is not composed purely of those exact primitive moves
accumulates a small rounding error at every step. The mean difference is therefore small
(most pixels are reasonably close to an edge, where few steps are needed and the
approximation is tight), while the max difference is larger, consistent with the
approximation error growing for pixels further from the nearest edge or along directions
less well represented by the three primitive move types.

\subsection{Chamfer Distance Between Two Fixed Edge Maps (Part 5b)}

\begin{center}
\begin{tabular}{lc}
\toprule
Direction & Value \\
\midrule
$d_{\text{ch}}(A,B)$ & $0.0000$ \\
$d_{\text{ch}}(B,A)$ & $1.2500$ \\
\bottomrule
\end{tabular}
\end{center}

The synthetic edge maps used: $A$ is a small L-shaped set of edge pixels on a $10\times10$
grid, and $B$ is exactly $A$ plus one additional isolated edge pixel elsewhere on the grid.

The two directional Chamfer distances $d_{\text{ch}}(A,B)$ and $d_{\text{ch}}(B,A)$ need not
be equal because each one averages, over a different point set ($A$ or $B$), the distance to
the \emph{other} set's nearest point --- and these are generally different quantities. In
particular, since $B$ is constructed to contain every point of $A$ plus at least one extra
point, every point of $A$ has distance $0$ to its nearest neighbour in $B$ (because that
neighbour can always be the corresponding point of $A$ itself, which is also in $B$), giving
$d_{\text{ch}}(A,B) = 0$. The reverse direction, $d_{\text{ch}}(B,A)$, includes the extra
point(s) in $B$ that are not in $A$; those contribute a strictly positive distance to their
nearest neighbour in $A$, so $d_{\text{ch}}(B,A) > 0$. This asymmetry is exactly why Chamfer
distance is a one-sided (directional) similarity measure rather than a true metric distance
between sets. This matches the measured values exactly: $d_{\text{ch}}(A,B) = 0$, since every
point of $A$ is also in $B$ and so has a self-match at distance $0$, while
$d_{\text{ch}}(B,A) = 1.25 > 0$, reflecting the one extra point in $B$ that has no exact match
in $A$ and must average in a non-zero distance to its nearest neighbour in $A$.

\end{document}
